\begin{frame}{보고서 구조 (M4) — 6절}
  8.1의 6절 구조를 그대로 쓰되, \textbf{2절과 5절이 Block B용으로
  바뀜}.
  \begin{enumerate}
    \item \kb{문제 정의}(0.5p) — 행의 물리적 단위, 라벨, 주 지표,
      \textbf{두 층의 베이스라인(무작위·GBDT)} + \textbf{정형/비정형
      판정과 그 근거}(축 1)
    \item \kb{데이터와 전처리}(1p) — 원본 통계, \textbf{train으로만
      fit한 전처리} 명시(이미지 0~1 스케일, 텍스트 토크나이저/패딩),
      분할 비율 + stratify/패딩 여부
    \item \kb{모델 선택과 이유}(1~2p) — 적용한 \textbf{딥러닝 기법
      $\geq$1개}와 그 구조(층·크기)를 Ch10.1 공식으로 정당화
      (§``손으로 한 번'') + \textbf{왜 GBDT가 아니라} 이 기법인가(축 2)
    \item \kb{결과}(1p) — 최종 test 숫자(한 번), \textbf{에포크 학습
      곡선 1장}, \textbf{GBDT·무작위와의 비교 표}
    \item \kb{수식적 정당화}(1p) — \textbf{필수 1개}(10.1 파라미터 /
      9.3 정규화 / 15 ELBO) — §``손으로 한 번''의 표를 그대로
    \item \kb{잘 안 된 부분과 원인 진단}(0.5p) — \textbf{필수}.
      학습이 안 된/과적합된 부분을 Ch09·09.3·10.2의 개념으로 진단 —
      8.1의 ``\textbf{모두 잘 됐다}'' 경고가 그대로 적용
  \end{enumerate}
\end{frame}
